% GENERATED FILE. Edit canonical lessons, then run npm run generate:books.
% canonical-source-hash: fnv1a64:45811306c2f50669
% canonical-lessons: TA-C318-tirappuvila, TA-C318-kuttuvilakku, TA-C318-cirappu-viruntinar, TA-C318-talaimaitanku, TA-C318-uraiyarru

\chapter{An Opening Ceremony}
\label{ch:ta-an-opening-ceremony}
\hlchaptermodality{318}

\begin{chapteropening}
\textbf{By the end of this chapter:} \emph{I can describe the opening of a building or a festival.}

Complete the last lesson of chapter 318: I can describe the opening of a building or a festival.
\end{chapteropening}

\section[\texorpdfstring{\ta{திறப்புவிழா} (tiṟappuviḻā)}{திறப்புவிழா (tiṟappuviḻā)}]{\ta{திறப்புவிழா} (tiṟappuviḻā) — an opening ceremony}

\label{lesson:TA-C318-tirappuvila}



\begin{quote}
Before the new one: say the Tamil for then and there, then the Tamil for one after another.
\end{quote}

\subsection*{You'll want to know: \ta{திறப்புவிழா}}
\textbf{\ta{திறப்புவிழா}} — \emph{tiṟappuviḻā} — \textquotedblleft{}an opening ceremony\textquotedblright{}.

\begin{quote}
\ta{புதிய} \ta{நூலகத்தின்} \ta{திறப்புவிழா} \ta{இன்று} \ta{நடந்தது}.
\end{quote}

\begin{quote}
\emph{putiya nūlakattiṉ tiṟappuviḻā iṉṟu naṭantatu.}
\end{quote}

\begin{quote}
The opening ceremony of the new library took place today.
\end{quote}

\ta{திறப்பு}, opening + \ta{விழா}, festival.

\begin{cousinweb}
\textbf{\ta{திறப்புவிழா}} (\emph{tiṟappuviḻā}) is a Tamil compound built from \ta{திறப்பு}, \textquotedblleft{}opening\textquotedblright{}, from \ta{திற}, \textquotedblleft{}to open\textquotedblright{} (DEDR 3259), and \ta{விழா}, \textquotedblleft{}festival\textquotedblright{} (DEDR 5427).
\end{cousinweb}

\begin{grammarlens}[title={What you've built}]
One more word for telling what happened.
\end{grammarlens}

\subsection*{Your turn}
\begin{itemize}
\raggedright
  \item \emph{Say it:} \emph{tiṟappuviḻā}
  \item \emph{Say it:} \emph{tiṟappuviḻā}, once more
  \item \emph{Say it:} \emph{putiya nūlakattiṉ tiṟappuviḻā iṉṟu naṭantatu.} — The opening ceremony of the new library took place today.
  \item \emph{Recall:} say the Tamil for then and there, then the Tamil for one after another, then say \emph{tiṟappuviḻā} again
\end{itemize}

\begin{tcolorbox}[breakable,colback=teal!4,colframe=teal!35!black,title={Before you move on}]
What does \ta{திறப்புவிழா} mean? (\textquotedblleft{}An opening ceremony\textquotedblright{}.) Say it once more.
\end{tcolorbox}
\section[\texorpdfstring{\ta{குத்துவிளக்கு} (kuttuviḷakku)}{குத்துவிளக்கு (kuttuviḷakku)}]{\ta{குத்துவிளக்கு} (kuttuviḷakku) — a standing lamp}

\label{lesson:TA-C318-kuttuvilakku}



\begin{quote}
Before the new one: say the Tamil for one after another, then the Tamil for an opening ceremony.
\end{quote}

\subsection*{You'll want to know: \ta{குத்துவிளக்கு}}
\textbf{\ta{குத்துவிளக்கு}} — \emph{kuttuviḷakku} — \textquotedblleft{}a standing lamp\textquotedblright{}.

\begin{quote}
\ta{விழா} \ta{குத்துவிளக்கு} \ta{ஏற்றித்} \ta{தொடங்கியது}.
\end{quote}

\begin{quote}
\emph{viḻā kuttuviḷakku ēṟṟit toṭaṅkiyatu.}
\end{quote}

\begin{quote}
The ceremony began with the lighting of the standing lamp.
\end{quote}

\ta{குத்துவிளக்கு} \ta{ஏற்று}, to light the lamp, opens many functions.

\begin{cousinweb}
\textbf{\ta{குத்துவிளக்கு}} (\emph{kuttuviḷakku}) is a Tamil compound built from \ta{குத்து}, \textquotedblleft{}to set upright, fix in the ground\textquotedblright{} (DEDR 1720), and \ta{விளக்கு}, \textquotedblleft{}a lamp\textquotedblright{}.
\end{cousinweb}

\begin{grammarlens}[title={What you've built}]
One more word for telling what happened.
\end{grammarlens}

\subsection*{Your turn}
\begin{itemize}
\raggedright
  \item \emph{Say it:} \emph{kuttuviḷakku}
  \item \emph{Say it:} \emph{kuttuviḷakku}, once more
  \item \emph{Say it:} \emph{viḻā kuttuviḷakku ēṟṟit toṭaṅkiyatu.} — The ceremony began with the lighting of the standing lamp.
  \item \emph{Recall:} say the Tamil for one after another, then the Tamil for an opening ceremony, then say \emph{kuttuviḷakku} again
\end{itemize}

\begin{tcolorbox}[breakable,colback=teal!4,colframe=teal!35!black,title={Before you move on}]
What does \ta{குத்துவிளக்கு} mean? (\textquotedblleft{}A standing lamp\textquotedblright{}.) Say it once more.
\end{tcolorbox}
\section[\texorpdfstring{\ta{சிறப்பு} \ta{விருந்தினர்} (ciṟappu viruntiṉar)}{சிறப்பு விருந்தினர் (ciṟappu viruntiṉar)}]{\ta{சிறப்பு} \ta{விருந்தினர்} (ciṟappu viruntiṉar) — a chief guest}

\label{lesson:TA-C318-cirappu-viruntinar}



\begin{quote}
Before the new one: say the Tamil for an opening ceremony, then the Tamil for a standing lamp.
\end{quote}

\subsection*{You'll want to know: \ta{சிறப்பு} \ta{விருந்தினர்}}
\textbf{\ta{சிறப்பு} \ta{விருந்தினர்}} — \emph{ciṟappu viruntiṉar} — \textquotedblleft{}a chief guest\textquotedblright{}.

\begin{quote}
\ta{விழாவுக்குச்} \ta{சிறப்பு} \ta{விருந்தினராக} \ta{ஒரு} \ta{மருத்துவர்} \ta{வந்தார்}.
\end{quote}

\begin{quote}
\emph{viḻāvukkuc ciṟappu viruntiṉarāka oru maruttuvar vantār.}
\end{quote}

\begin{quote}
A doctor came to the ceremony as the chief guest.
\end{quote}

-\ta{ஆக}, as: \ta{சிறப்பு} \ta{விருந்தினராக}, as chief guest.

\begin{cousinweb}
\textbf{\ta{சிறப்பு} \ta{விருந்தினர்}} (\emph{ciṟappu viruntiṉar}) joins native \ta{சிறப்பு}, \textquotedblleft{}pre-eminence, honour\textquotedblright{} (DEDR 2589), and \ta{விருந்தினர்}, \textquotedblleft{}a guest\textquotedblright{} (DEDR 5415).
\end{cousinweb}

\begin{grammarlens}[title={What you've built}]
One more word for telling what happened.
\end{grammarlens}

\subsection*{Your turn}
\begin{itemize}
\raggedright
  \item \emph{Say it:} \emph{ciṟappu viruntiṉar}
  \item \emph{Say it:} \emph{ciṟappu viruntiṉar}, once more
  \item \emph{Say it:} \emph{viḻāvukkuc ciṟappu viruntiṉarāka oru maruttuvar vantār.} — A doctor came to the ceremony as the chief guest.
  \item \emph{Recall:} say the Tamil for an opening ceremony, then the Tamil for a standing lamp, then say \emph{ciṟappu viruntiṉar} again
\end{itemize}

\begin{tcolorbox}[breakable,colback=teal!4,colframe=teal!35!black,title={Before you move on}]
What does \ta{சிறப்பு} \ta{விருந்தினர்} mean? (\textquotedblleft{}A chief guest\textquotedblright{}.) Say it once more.
\end{tcolorbox}
\section[\texorpdfstring{\ta{தலைமைதாங்கு} (talaimaitāṅku)}{தலைமைதாங்கு (talaimaitāṅku)}]{\ta{தலைமைதாங்கு} (talaimaitāṅku) — to preside}

\label{lesson:TA-C318-talaimaitanku}



\begin{quote}
Before the new one: say the Tamil for a standing lamp, then the Tamil for a chief guest.
\end{quote}

\subsection*{You'll want to know: \ta{தலைமைதாங்கு}}
\textbf{\ta{தலைமைதாங்கு}} — \emph{talaimaitāṅku} — \textquotedblleft{}to preside\textquotedblright{}.

\begin{quote}
\ta{ஊர்ப்} \ta{பெரியவர்} \ta{விழாவுக்குத்} \ta{தலைமைதாங்கினார்}.
\end{quote}

\begin{quote}
\emph{ūrp periyavar viḻāvukkut talaimaitāṅkiṉār.}
\end{quote}

\begin{quote}
A village elder presided over the ceremony.
\end{quote}

Past (I …): \ta{தலைமைதாங்கினேன்}, \emph{talaimaitāṅkiṉēṉ}.

The event takes -\ta{க்கு}: \ta{விழாவுக்குத்} \ta{தலைமைதாங்கு}.

\begin{cousinweb}
\textbf{\ta{தலைமைதாங்கு}} (\emph{talaimaitāṅku}) is a Tamil compound built from \ta{தலைமை}, \textquotedblleft{}leadership\textquotedblright{} (DEDR 3103), and \ta{தாங்கு}, \textquotedblleft{}to bear, support\textquotedblright{} (DEDR 3153): to bear the lead.
\end{cousinweb}

\begin{grammarlens}[title={What you've built}]
One more word for telling what happened.
\end{grammarlens}

\subsection*{Your turn}
\begin{itemize}
\raggedright
  \item \emph{Say it:} \emph{talaimaitāṅku}
  \item \emph{Say it:} \emph{talaimaitāṅku}, once more
  \item \emph{Say it:} \emph{ūrp periyavar viḻāvukkut talaimaitāṅkiṉār.} — A village elder presided over the ceremony.
  \item \emph{Recall:} say the Tamil for a standing lamp, then the Tamil for a chief guest, then say \emph{talaimaitāṅku} again
\end{itemize}

\begin{tcolorbox}[breakable,colback=teal!4,colframe=teal!35!black,title={Before you move on}]
What does \ta{தலைமைதாங்கு} mean? (\textquotedblleft{}To preside\textquotedblright{}.) Say it once more.
\end{tcolorbox}
\section[\texorpdfstring{\ta{உரையாற்று} (uraiyāṟṟu)}{உரையாற்று (uraiyāṟṟu)}]{\ta{உரையாற்று} (uraiyāṟṟu) — to give a speech}

\label{lesson:TA-C318-uraiyarru}



\begin{quote}
Before the new one: say the Tamil for a chief guest, then the Tamil for to preside.
\end{quote}

\subsection*{You'll want to know: \ta{உரையாற்று}}
\textbf{\ta{உரையாற்று}} — \emph{uraiyāṟṟu} — \textquotedblleft{}to give a speech\textquotedblright{}.

\begin{quote}
\ta{சிறப்பு} \ta{விருந்தினர்} \ta{மாணவர்களிடம்} \ta{உரையாற்றினார்}.
\end{quote}

\begin{quote}
\emph{ciṟappu viruntiṉar māṇavarkaḷiṭam uraiyāṟṟiṉār.}
\end{quote}

\begin{quote}
The chief guest addressed the students.
\end{quote}

Past (I …): \ta{உரையாற்றினேன்}, \emph{uraiyāṟṟiṉēṉ}.

Those addressed take -\ta{இடம்}; a formal word for \ta{பேசு}.

\begin{cousinweb}
\textbf{\ta{உரையாற்று}} (\emph{uraiyāṟṟu}) is a Tamil compound built from \ta{உரை}, \textquotedblleft{}speech\textquotedblright{} (DEDR 648), and \ta{ஆற்று}, \textquotedblleft{}to accomplish, perform\textquotedblright{} (DEDR 407).
\end{cousinweb}

\begin{grammarlens}[title={What you've built}]
One more word for telling what happened.
\end{grammarlens}

\subsection*{Your turn}
\begin{itemize}
\raggedright
  \item \emph{Say it:} \emph{uraiyāṟṟu}
  \item \emph{Say it:} \emph{uraiyāṟṟu}, once more
  \item \emph{Say it:} \emph{ciṟappu viruntiṉar māṇavarkaḷiṭam uraiyāṟṟiṉār.} — The chief guest addressed the students.
  \item \emph{Recall:} say the Tamil for a chief guest, then the Tamil for to preside, then say \emph{uraiyāṟṟu} again
\end{itemize}

\begin{tcolorbox}[breakable,colback=teal!4,colframe=teal!35!black,title={Before you move on}]
What does \ta{உரையாற்று} mean? (\textquotedblleft{}To give a speech\textquotedblright{}.) Say it once more.
\end{tcolorbox}
